\documentclass[tikz,border=5pt]{standalone}
\usetikzlibrary{arrows.meta}
\usetikzlibrary{angles,quotes,calc}

\begin{document}
	\begin{tikzpicture}[>=Latex]
		% 坐标轴
		\draw[->,thick] (0,0) -- (18,0) node[below] {X(East)};
		\draw[->,thick] (0,0) -- (0,12) node[left] {Y(North)};
		
		% 两个航点
		\coordinate (Pk) at (2, 5);
		\coordinate (Pkp) at (16, 13.17);
		
		% 路径线
		\draw[thick] (Pk) -- (Pkp);
		
		% 航点
		\fill (Pk) circle (1.2pt);
		\fill (Pkp) circle (1.2pt);
		
		% 标签
		\node[below left] at (Pk) {$p_k$};
		\node[above right] at (Pkp) {$p_{k+1}$};
		
		% 船的位置
		\coordinate (P) at (7.5, 1);
		\fill (P) circle (1.2pt);
		\node[above left] at (P) {$p$};
		
		
		% 投影点（手工给定）
		\coordinate (Pproj) at (4.5, 6.46);
		\fill (Pproj) circle (1.2pt);
		\node[below] at (Pproj) {$p_\perp$};
		
		% 从船到投影点的虚线   % 横向误差线
%		\draw[dashed] (P) -- (Pproj);
		\draw[thick] (P) -- (Pproj);
		\node[right] at (5, 4) {$e$};
		
		  % LOS 点
		\coordinate (Plos) at (14, 12);
		\fill (Plos) circle (1.2pt);
		\node[above] at (Plos) {$p_{\mathrm{LOS}}$};
		
		% 前视距离	
		\draw[<->] (Pproj) -- (Plos);
		\node[below] at (10, 10.5) {$\Delta$};
		
		% LOS 向量
		\draw[->,thick] (P) -- (Plos);
		
		\coordinate (P1) at (7.5, 10);
		\draw[dashed] (P) -- (P1);
		
		% 用于角度标注的 East 参考点
		\coordinate (Peast) at ($(P)+(0,7.21)$);	
		
		% alpha_k
		\draw pic[
		draw,
		<-,
		angle radius=9mm,
		angle eccentricity=1.25,
		"$\alpha_k$"
		] {angle = Pkp--Peast--P1};
		
		\coordinate (P2) at (14, 10);
		\draw[dashed] (Plos) -- (P2);
		
		% chi_d 角
%		\coordinate (PeastP) at ($(P)+(1.4,0)$);
		\draw pic[
		draw,
		<-,
		angle radius=8mm,
		angle eccentricity=1.25,
		"$\chi_d$"
		] {angle = P--Plos--P2};
		

		% -------------------------------------------------
		% 3. 定义三个方向角
		%    chi_d: 期望航迹角（LOS方向）
		%    chi  : 实际航迹角
		%    psi  : 实际航向角（船体方向）
		% -------------------------------------------------
		% 计算chiD 
		\pgfmathanglebetweenpoints{\pgfpointanchor{P}{center}}{\pgfpointanchor{Plos}{center}}
		\let\chiD\pgfmathresult
		
		% 你可以手动调这两个偏差，控制图的观感
		\pgfmathsetmacro{\chiActual}{\chiD+15}   % 实际航迹角 chi
		\pgfmathsetmacro{\psiAngle}{\chiActual-10} % 实际航向角 psi
		% beta = chi - psi = 16 deg

		% -------------------------------------------------
		% 4. 先定义三条方向辅助点
		% -------------------------------------------------
		\coordinate (ChiDPoint) at ($(P)+({4*cos(\chiD)},{4*sin(\chiD)})$);
		\coordinate (ChiPoint)  at ($(P)+({3*cos(\chiActual)},{3*sin(\chiActual)})$);
		\draw[->, dashed] (P) -- (ChiPoint) node[above] {$U$};
		\coordinate (PsiPoint)  at ($(P)+({3*cos(\psiAngle)},{3*sin(\psiAngle)})$);
%		\draw[dashed] (P) -- (PsiPoint);
		\coordinate (NorthRef) at $(P)+(0, 1.5);
		
		% =========================
		% 3. 船体参数（从你原来的代码缩小后搬过来）
		% =========================
		\def\L{1.6}   % 船全长
		\def\B{0.6}   % 船身宽
%		\def\Bbow{0.2} % 船首尖端宽 (0表示全尖)
		\def\Lrect{1.0} % 矩形部分长度
%		\def\psiAngle{35} % USV 的航向角
		% =========================
		% 4. 在 P 点画 USV
		% =========================
		\begin{scope}[shift={(P)}, rotate=\psiAngle]
			
			\coordinate (SternL) at (-\Lrect/2, -\B/2);
			\coordinate (SternR) at (-\Lrect/2,  \B/2);
			\coordinate (BowR)   at ( \Lrect/2,  \B/2);
			\coordinate (BowTip) at ( \L/2,      0);
			\coordinate (BowL)   at ( \Lrect/2, -\B/2);
			
%			\draw[thick, =blue!10, fill opacity=0.25]
			\draw[thick]
			(SternL) -- (SternR) -- (BowR) -- (BowTip) -- (BowL) -- cycle;
			
			\draw[->, dashed] %, red!80!black]
			(0,0) -- (3,0) coordinate (Xb_axis) node[above] {$x_B$};
			
			\draw[->, dashed] %, red!80!black]
			(0,0) -- (0,-1) node[above] {$y_B$};
			
		\end{scope}
		
		% -------------------------------------------------
		% 7. 角度标注：psi, chi, beta
		% -------------------------------------------------
		
		% psi：从 East 到船体航向
		\draw pic[
		draw,
		<-,
		angle radius=25mm,
		angle eccentricity=1.1,
		"$\psi$"
		] {angle = PsiPoint--P--NorthRef};
		
		% chi：从 East 到实际航迹
		\draw pic[
		draw,
		<-,
		angle radius=18mm,
		angle eccentricity=1.18,
		"$\chi$"
		] {angle = ChiPoint--P--NorthRef};
		
		% beta：从 psi 到 chi
		\draw pic[
		draw,
		<-,
		angle radius=19mm,
		angle eccentricity=1.18,
		"$\beta$"
		] {angle = PsiPoint--P--ChiPoint};
		
		% chi_d：期望航迹
		\draw pic[
		draw,
		<-,
		angle radius=12mm,
		angle eccentricity=1.18,
		"$\chi_d$"
		] {angle = ChiDPoint--P--NorthRef};

	\end{tikzpicture}
\end{document}